\section{The Rodden Sorting Constraint}
\label{sec:constraints}

\subsection{Geographic Sorting as a Structural Binding Constraint}

The most important single finding from the E-track experiments is that the
Rodden geographic-sorting constraint \citep{rodden2019why} operates differently
across system types and is the binding constraint that separates systems into
two fundamentally different classes.

The sorting constraint arises from the empirical fact that Democratic voters are
geographically concentrated in urban cores at densities substantially higher than
Republican voters anywhere.  \citet{chen2013unintentional} show that this
concentration produces Republican seat bonuses under any compact single-member
district drawing, algorithmic or otherwise.  \citet{rodden2019why} quantifies
the mechanism: Democrats' geographic inefficiency (many votes wasted in
landslide urban wins) translates to proportional underrepresentation even
under neutral drawing procedures.

The implications for system design depend critically on whether the system
requires geographic districts:

\begin{enumerate}
  \item \textbf{Single-member geographic systems} (Baseline DIA, E.3, E.4, E.5):
        The sorting constraint applies fully.  Proportional outcomes are
        geometrically infeasible because drawing compact districts in areas with
        dense Democratic concentrations necessarily produces landslide Democratic
        wins that waste votes.

  \item \textbf{Multi-member geographic systems} (E.1):
        The sorting constraint applies partially.  Larger districts contain more
        geographic diversity, reducing the wasted-vote problem.  The degree of
        escape depends on district size: three-member districts reduce the
        sorting gap by approximately 40\%; five-member districts reduce it by
        approximately 65\% \citep{dellaLibera2026e1}.  Neither fully escapes
        sorting because urban cores remain more Democratic than any realistically
        sized district can accommodate proportionally.

  \item \textbf{Non-geographic systems} (E.2 county, E.5 party-based):
        The sorting constraint does not apply in the same form, because seats
        are allocated by formula rather than drawn geographically.  E.2 replaces
        the sorting problem with a county-size inequality problem: large
        multi-member counties (LA County, Cook County) elect at-large
        representatives who face effectively proportional elections, while small
        single-seat rural counties have winner-take-all outcomes.

\end{enumerate}

\subsection{How Single-Member Systems Are Affected}

All single-member systems --- including the algorithmic baseline --- are subject
to the sorting constraint in essentially the same form.  The key result from E.3
is instructive: even when state boundaries are removed and the algorithm
optimizes compactness globally, the proportionality outcome does not improve
significantly.  This is because the sorting constraint is not a consequence of
state boundaries or drawing-method idiosyncrasies --- it is a consequence of
where Democratic voters live relative to Republican voters \citep{rodden2019why}.

Similarly, E.4 (partisan-similarity clustering) demonstrates that changing the
edge-weighting scheme can change the \emph{distribution} of partisan margins
--- creating more lopsided districts --- but cannot change the aggregate
seat-vote relationship when total seat counts are fixed by population.

E.5 (partisan-fairness allocation) is the most direct confrontation with the
sorting constraint.  To achieve proportionality, the algorithm must draw
certain urban districts as geographic outliers --- non-compact, elongated, or
fragmented --- in order to reduce the Democratic wasted-vote surplus.  The
resulting maps are technically compliant with population equality and
contiguity requirements but would fail any serious compactness criterion.

\subsection{How Multi-Member Systems Partially Escape}

The partial escape of multi-member districts from the sorting constraint is
the central empirical finding of E.1.  The mechanism is district-size
averaging: a three-member district centered on an urban core will have a
population that includes suburban rings as well as the urban center, reducing
the Democrat-concentration effect that drives wasted votes.

The 40\% gap reduction for three-member districts and 65\% reduction for
five-member districts are derived from simulations across all 50 states in
which we compare the seat-vote deviation for single-member, three-member, and
five-member district configurations holding the geographic algorithm constant
\citep{dellaLibera2026e1}.  The residual gap --- the 35\% and 60\% of sorting
that remains even with five-member districts --- reflects the fact that some
urban concentrations are so extreme (Manhattan, central Chicago, central Los
Angeles) that no realistic district boundary can make them proportional.

\subsection{County Systems and the Sorting Bypass}

County-based representation (E.2) eliminates sorting in competitive states but
reintroduces it in a different form in highly urbanized states.  A county like
New York County (Manhattan) would receive approximately 2.3 House seats under
Huntington--Hill allocation and would elect those seats at-large by county vote
share.  Since Manhattan votes approximately 85--15 Democratic, the at-large
outcome would be 2 Democratic seats and 0.3 of a Republican seat (rounded to 0
in any whole-seat implementation).  This is more proportional than a single-
member district (which would produce 1--0 outcomes) but still not perfectly
proportional.

The deeper issue is that county boundaries do not respect partisan geography
any more than census-tract boundaries do.  They do, however, eliminate the
\emph{drawing} step that allows partisan manipulation, trading the gerrymandering
risk for the fixed-boundary risk that county demographics may themselves be
unrepresentative.
